\section{Attribution, dependencies, and verification}\label{app:attribution}

This draft was prepared with OpenAI Codex at the user's request.
The earlier drafts in the same chat supplied some of the expanded
elementary calculations. Their underlying mathematical sources remain
Fauzan, Lai, Pr\'evost, and the classical identities listed below.
Reusing those explanations does not make their mathematics original.

This revision also draws on the chat
\emph{Find parallel constructions in the Zeta(5)\ldots}
and its corrected explanatory companion.
They suggested making the common bilinear form, the roles of the local
bases, the passage from finite products and counts to continuous
estimates, and the primitive evaluation size explicit.
The chat's final provenance correction is retained: Fauzan already
states the underlying connections, including the weighted pullback,
the local-basis strategy, the periodic primitives, and the comparison
measure argument. The chat contributes organization and classical
interpretations. It supplies no new quantitative estimate.

\begin{longtable}{>{\raggedright\arraybackslash}p{0.27\textwidth}>{\raggedright\arraybackslash}p{0.65\textwidth}}
\toprule
Part of this account & Source and nature of the dependence\\
\midrule
\endhead
Guiding narrative &
Calegari's Persiflage post: the positive-measure, rational-modifier,
Hankel-determinant viewpoint and the connection with prior approximation
theory. It precedes Lai's exposition.\\[3pt]
Positive measure and transform &
Pr\'evost, Theorems 1--2. Euler's and Hermite's classical identities
are also recorded in DLMF. The special fifth-order calculation is
expanded here, as in Fauzan and Lai.\\[3pt]
Concrete determinant &
Fauzan, Section 2: \(D_N^6/D_K\), \(K=40n,N=3n\), affine entries,
degree, positivity, and the factorial scalar. Lai is a cross-check
and an exposition used in the preceding drafts.\\[3pt]
General determinant context &
Zudilin's irrationality criterion; Brown's integral pairing and
normalization framework; Uvarov's rational modifications, as discussed
and extended by Krattenthaler. These are contextual antecedents,
not additional unproved estimate inputs.\\[3pt]
Analytic estimate &
Fauzan, Section 6 and Appendix A: external field, comparison measure,
configuration inequality, and interval partition. Firsching's analytic
files supply detailed cross-checks. The finite certificate was re-executed
for the detailed version; its checked files accompany this condensation.\\[3pt]
Local arithmetic &
Fauzan's weighted-basis strategy, in the revised finite
partial-fraction implementation of Firsching's \texttt{Arith} files.
The full rank loss and cutoff 400 are retained.\\[3pt]
Global arithmetic &
Firsching's \texttt{Norm} and \texttt{Growth} files: allocation,
prime ranges, continuous profiles and corrected numerical total.
The rational pieces and integrals were independently regenerated
for the detailed version; the same data are used here.\\[3pt]
Parallel-construction explanations &
The named chat and its revised companion: the common-pairing
presentation, Euler--Maclaurin language, periodic Bernoulli notation,
and scale-independent evaluation size. Their mathematical inputs
are credited to Fauzan, the classical Bernoulli identities, and
Firsching's revised estimates at the relevant steps.\\[3pt]
Finite computation &
The companion files identify which data are transcribed from Fauzan
or Firsching and which verification code was written for this draft.
Any adapted repository code retains its license and attribution.\\
\bottomrule
\end{longtable}

The local source ledger records URLs, versions, hashes, and precise
locations. The mathematical proof and its finite checks should be
distinguished from a complete local rebuild of the formal development:
no such rebuild is claimed. The source repository's CI results are
external evidence, not a substitute for the calculations described here.

The uniform potential bound is certified interval by interval, including
the range beyond the finite partition. The arithmetic tables are checked
with rational arithmetic on complete intervals. Neither certificate is
based on testing selected floating-point values.

The background theorems are Tonelli and Fubini, elementary Fourier
convergence, finite dimensional linear algebra, the prime number theorem,
and the classical Bernoulli and factorial facts. Argument-specific
supporting estimates cited in the concise sections are proved in the
separate detailed companion.
